\documentclass[aspectratio=169,11pt]{beamer}
\usetheme{default}
\usecolortheme{default}
\setbeamertemplate{navigation symbols}{}
\setbeamertemplate{footline}[frame number]
\setbeamerfont{footnote}{size=\scriptsize}
\usepackage{booktabs}
\usepackage{tabularx}

\input{macros}

\title{Email Recovery Without Plaintext}
\subtitle{Zero-Knowledge Email as a Recovery Factor for Smart Accounts}
\author{Anonymous Author(s)}
\date{}

\begin{document}

\begin{frame}
  \titlepage
  \centering\small
  Every claim carries its status:
  \claimtag{proven} \claimtag{tested} \claimtag{measured} \claimtag{assumed} \claimtag{unproven}\\[2pt]
  \footnotesize\claimtag{proven} means a Lean theorem about a model of the contracts and the circuit's relation
\end{frame}

\begin{frame}{The problem}
  \begin{itemize}
    \item A passkey smart account is lost with the last device that holds the passkey.
    \item An inbox is a credential people keep for years~\cite{bonneau-recovery}.
    \item Naive email recovery puts the address in a database and makes the
      checker a party that can recover any account~\cite{sep30}.
    \item Checking DKIM on chain removes the server but publishes the
      \texttt{From} header.
  \end{itemize}
  \bigskip
  \textbf{Thesis:} inbox control can be \emph{one} recovery factor while the
  chain holds only a salted commitment and the operator's servers see the
  address only inside an attested enclave. Email never skips the recovery
  delay; alone it authorizes only under email plus delay (7 days, a second
  alert channel), where the owner's veto is the defence.
\end{frame}

\begin{frame}{The approach in one picture}
  \centering
  \resizebox{0.92\linewidth}{!}{\input{figures/recovery-flow}}
  \par\medskip\small
  A per-account guardian with $\Cm = \commitf(\addr, \code)$, counted as one
  guardian in Perch's recovery, used unchanged\claim{guardian-on-perch}.
\end{frame}

\begin{frame}{What the user sees}
  \small
  \begin{columns}[T]
    \begin{column}{0.48\textwidth}
      \textbf{Enrollment} (the passkey still works)
      \begin{enumerate}\footnotesize
        \item Save a 26-character recovery code; type two groups back.
        \item Send one prepared email from the inbox.
        \item Give an alert address, ideally at another provider; allow
          notifications on this device.
        \item The wallet deploys and lists the email helper.
      \end{enumerate}
      \footnotesize\claim{enrollment-safeguards}
    \end{column}
    \begin{column}{0.48\textwidth}
      \textbf{Recovery} (every passkey device lost)
      \begin{enumerate}\footnotesize
        \item Open the recovery page from the kit's account address.
        \item Type the code under ``nobody will ask for this''.
        \item Send one prepared email; its body says to stop if anyone asked.
        \item Wait for the proof; other helpers approve; the delay passes.
      \end{enumerate}
      \footnotesize\claims{recovery-page-aids,phishing-defences}
    \end{column}
  \end{columns}
  \medskip
  A device that still holds the passkey is told it needs no recovery and
  offered the cancellation\claim{wallet-live-attempts}.
\end{frame}

\begin{frame}{Background}
  \begin{description}[Perch recovery]\small
    \item[DKIM] The signing domain \texttt{d=} signs the fields \texttt{h=}
      names and the signature field itself, under a key in DNS; relaxed
      canonicalization unfolds lines~\cite{rfc6376,rfc8301}. It attests to
      the domain, not the author.
    \item[Perch recovery] Anyone opens an attempt; guardians approve the
      SHA-256 digest of a fixed-width statement; at the quorum the attempt is
      authorized, and it completes after a delay during which a
      \texttt{Loss} owner can cancel~\cite{perch-recovery-spec}.
    \item[UltraHonk] Noir circuits proved by Barretenberg, verified on
      Soroban by Perch's verifier; proofs are 16{,}224 bytes\claim{circuit-size}.
  \end{description}
\end{frame}

\begin{frame}{Threat model}
  \footnotesize
  \begin{columns}[T]
    \begin{column}{0.42\textwidth}
      \textbf{Eight adversaries}
      \begin{itemize}
        \item inbox thief
        \item social engineer (chosen email, read-out code)
        \item registry members, below or at threshold
        \item relayer or front-runner
        \item mail provider (curious, not forging)
        \item network observer
        \item malicious frontend (trusted at enrollment)
        \item the operator's staff (host, DNS)
      \end{itemize}
    \end{column}
    \begin{column}{0.56\textwidth}
      \textbf{Assumptions}
      \begin{itemize}
        \item A1a DKIM unforgeable; A1b providers authenticate senders\claims{dkim-unforgeability,provider-authenticates-sender}
        \item A2, A3 hashes; A4 UltraHonk; A5 registry; A6 secrets; A7 enrollment\claims{hash-collision-resistance,proof-system-security,registry-correctness,secret-generation,honest-enrollment}
        \item A8 enclave and custodians\claim{enclave-isolation}
        \item A9 close time and ledger rate\claim{ledger-timing}
        \item A10 recovery domain's MX, MTA-STS and TLS keys unchanged\claim{recovery-domain-routing}
        \item A11 the published wallet build\claim{published-wallet-build}
        \item A12 under email plus delay, the owner learns of an attempt in time\claim{alert-liveness}
      \end{itemize}
    \end{column}
  \end{columns}
\end{frame}

\begin{frame}{Enrollment}
  \centering
  \input{figures/sequence-enroll}
\end{frame}

\begin{frame}{Recovery}
  \centering
  \input{figures/sequence-recover}
\end{frame}

\begin{frame}{Binding the statement, blinding it from the provider}
  \small
  \begin{itemize}
    \item The email approves Perch's own statement $S$; the guardian checks
      $\Sha(S) = \digest$ and that $S$ names its account. Account, new key,
      nonce and expiry are fields of $S$.
    \item Subject: \texttt{APPROVE \OpNameUpper{} ACCOUNT RECOVERY}
      $\mathrm{hex}(T)$, $T = \blindf(\digest, \blind)$, $\blind$ kept in the
      browser\claim{subject-blinding}.
    \item $T$'s content names no attempt\claim{blinded-subject-hiding}; but
      the vote lands seconds after the email, so a provider watching the
      chain can link them by timing, possibly to a set of
      one\claims{browser-proving-cost,testnet-email-recovery}.
    \item The mail reaches the browser sealed by the vault, or as an
      \texttt{.eml} file\claim{sealed-inbound}.
  \end{itemize}
\end{frame}

\begin{frame}{The circuit relation}
  \centering
  \resizebox{!}{0.78\textheight}{\input{figures/circuit-relation}}
  \par\smallskip\footnotesize
  248{,}996 gates; headers up to 1536 bytes\claims{circuit-size,header-bound-fits}.
  Only the signed bytes; a \texttt{From} an RFC~5322 reader reads as exactly
  the bound mailbox; aligned \texttt{d=}; any public input changed, the proof
  fails\claims{circuit-from-grammar,circuit-dkim-rules,circuit-binds-public-inputs}.
\end{frame}

\begin{frame}{The email guardian}
  \small
  \begin{itemize}
    \item One immutable guardian per account: its address is in Perch's
      \texttt{config\_hash}, so changing the email factor passes Perch's
      reconfiguration gate.
    \item Twelve checks in order: one controller call with $\digest$;
      $\Sha(S)=\digest$; network, account, controller; allowed action and
      entry point; not expired; root is the live pin or registry-accepted;
      the proof verifies\claim{guardian-checks}.
    \item The verdict ignores the authorization entry\claim{guardian-payload-irrelevant};
      a trapping registry or verifier is a refusal.
    \item Construction refuses bad commitments, pins over 365 days, and a
      verifier for another circuit\claim{guardian-construction}.
  \end{itemize}
\end{frame}

\begin{frame}{DKIM key registry and the per-account pin}
  \small
  \begin{itemize}
    \item Additions: $k$-of-$n$ approval, then 72 hours of close time if
      attested live in DNS, else 7 days; allowlisted domains, RSA-2048
      only\claims{registry-timelocks,registry-governance}.
    \item Revocation refuses every root holding the key at once; no admin or
      upgrade\claims{registry-revocation,registry-immutable}.
    \item Pin: a salted key tree fixed at construction, additive, at most 365
      days of close time; the wallet uses 30\claims{key-pin,wallet-pin-lifetime}.
    \item A planted, activated key still cannot approve without the account
      code\claim{registry-timelock-attacks}.
  \end{itemize}
\end{frame}

\begin{frame}{Alerts, and the privacy claim's conditions}
  \small
  \begin{itemize}
    \item From Perch's events and state: an authorized attempt, and an email
      vote below quorum, on every channel; opening an attempt alerts
      nothing\claims{alert-code,compromised-inbox-alerted}.
    \item Web Push to a passkey device required, probed weekly; a lapse is
      shown and blocks keeping the email helper\claim{wallet-second-alert-channel}.
    \item Vault key 2-of-3; enclave images after a 14-day public
      timelock\claim{vault-custody}.
    \item Mocked alert path: no server-side record holds the
      address\claim{servers-cannot-read-address}.
    \item The wallet's claim, \PrivacyClaim, holds only under A8, A10,
      A11\claim{scoped-privacy-claim}. The mail provider and the sending
      service see the address.
  \end{itemize}
\end{frame}

\begin{frame}{Beyond the protocol: the human}
  \small
  \textbf{Social engineer:} the real recovery page on its own device opens
  the attempt and shows the Subject; the victim's email is a valid approval,
  and the attacker sees the same alert code.\par\medskip
  \textbf{What remains}
  \begin{itemize}
    \item the delay and the owner's cancellation; the push leads with
      cancelling; the wallet reads open attempts itself\claims{phishing-defences,wallet-live-attempts}
    \item other guardians; warnings where the email is prepared and the code typed
  \end{itemize}
  \textbf{Residual risk under email plus delay}
  \begin{itemize}
    \item a phished user with no device left cannot cancel
    \item alerts to a stolen inbox and a lapsed push: only opening the wallet
      shows the attempt
    \item whether people act in time is not measured\claim{user-behaviour}
  \end{itemize}
\end{frame}

\begin{frame}{What is proven, tested and assumed}
  \small
  \begin{description}[Soundness]
    \item[Soundness] An accepted approval yields, for every extractor, the
      email facts or the extractor's failure on that proof; with A5 and A7,
      the enrolled address and a trusted key, or a Poseidon2 collision on
      inputs the theorem computes\claim{email-factor-soundness}.
    \item[Privacy] In an ideal model (handles for $\Cm$ and the pin root,
      simulated proofs), an adaptive adversary's view, proofs included, is
      simulated from public data\claim{address-privacy}; the step to
      Poseidon2 and UltraHonk is assumed\claims{commitment-hiding,proof-system-security}.
  \end{description}
  \medskip
  Lean~4, no axioms beyond Lean's own\claim{lean-axioms}. Not proven: the
  circuit's parsing and the model's agreement with the code, both
  tested\claim{model-code-agreement}.
\end{frame}

\begin{frame}{Properties: replay and one-factor safety}
  \small
  \begin{description}[One factor]
    \item[Replay] An approval is accepted only for its own digest, which fixes
      one account, network, controller, action, attempt and epoch, or two
      named encodings collide; the controller counts the guardian
      once\claims{replay-resistance,digest-fixes-statement}.
    \item[One factor] With email approvals as the only evidence, an attempt is
      authorized only at quorum one, guardian-only, and completes only after
      a positive delay\claim{one-factor-safety}.
    \item[Wallet] Email alone reaches the quorum only under \texttt{Loss},
      with 7 days at the target close time and a second alert
      channel\claim{wallet-policy-one-factor}; in wall-clock time under A9.
  \end{description}
\end{frame}

\begin{frame}{The circuit's parsing}
  \footnotesize
  \begin{tabularx}{\linewidth}{@{}>{\raggedright\arraybackslash}p{0.36\linewidth}>{\raggedright\arraybackslash}X@{}}
    \toprule
    Public issue class & Refusal \\
    \midrule
    Display names spoofing \texttt{From} & one angle-address ending the field \\
    Folded \texttt{From}; simple canonicalization & no folded line; relaxed only \\
    Fields spanning lines; parser resets & a field starts at byte 0 or after CRLF \\
    $\sigma + N$; \texttt{bh=}, \texttt{l=}, misaligned \texttt{d=} & $\sigma < N$; scanned tags; aligned \texttt{d=} \\
    Unsigned or duplicated headers & only the signed block; one of each \\
    \bottomrule
  \end{tabularx}
  \par\smallskip
  One negative test per class~\cite{audit-matterlabs-2024,audit-veridise-2024,audit-consensys-2024}\claims{circuit-issue-classes,circuit-refuses-attacks}.
  \par\medskip
  \small
  \begin{itemize}
    \item Real mail: 5{,}327 signatures from all 11 launch domains, the same
      verdict as mailauth and Python on every one\claim{parser-real-mail}.
    \item Fuzzing: 50{,}000 headers; every accepted \texttt{From} is the
      mailbox Python reads\claim{parser-fuzzing}.
  \end{itemize}
\end{frame}

\begin{frame}{Evaluation: attack suite}
  \small
  \begin{columns}[T]
    \begin{column}{0.5\textwidth}
      \footnotesize
      \begin{tabular}{@{}lrr@{}}
        \toprule
        Class & Rows & Refused \\
        \midrule
        Forged emails & 11 & 10 \\
        Replays & 12 & 12 \\
        Wrong account or key & 5 & 5 \\
        Registry keys & 8 & 7 \\
        Compromised inbox & 12 & 7 \\
        Alert path & 3 & --- \\
        Phishing & 6 & 2 \\
        Repeated attempts & 1 & --- \\
        Header quirks & 5 & 2 \\
        Front-running & 4 & 4 \\
        Privacy on chain & 1 & --- \\
        \bottomrule
      \end{tabular}
    \end{column}
    \begin{column}{0.48\textwidth}
      \footnotesize
      68 rows, each with its expected outcome: 65 on a local network, 3 on
      the test host\claims{network-attack-suite,registry-timelock-attacks}.\\[4pt]
      Thief with the code, email plus delay: deletes the alert emails; the
      push reaches the owner, who cancels; a missed event is found by
      reconciliation\claim{compromised-inbox-alerted}.\\[4pt]
      Residual: a phished user with no device; a lapsed push.\\[4pt]
      No address or code in 213 transactions\claim{privacy-no-address-on-chain}.
    \end{column}
  \end{columns}
\end{frame}

\begin{frame}{Evaluation: what an email vote costs}
  \small
  \begin{tabular}{@{}lrrr@{}}
    \toprule
    & Test host & Local network & Testnet \\
    \midrule
    Settings & SDK snapshot & mainnet's & testnet's \\
    Instructions, metered (M) & 131.2 & 121.6 & 121.6 \\
    Transaction (bytes) & 18{,}148 & 18{,}080 & 18{,}080 \\
    Fee (stroops) & 178{,}196 & 83{,}125--83{,}155 & 82{,}424--82{,}442 \\
    \bottomrule
  \end{tabular}
  \par\smallskip{\footnotesize\claims{onchain-verification-cost,network-verification-cost,testnet-email-recovery}}
  \par\medskip
  \begin{itemize}\footnotesize
    \item Mainnet projection: 83{,}150 for the Gmail vote: 37{,}971 compute,
      36{,}419 history (an 18\,KB envelope), 2{,}049 rent\claim{fee-attribution}.
    \item The test host's SDK prices with an old fee snapshot; at mainnet's
      settings the same transaction is 82{,}530\claim{onchain-verification-cost}.
    \item Everything archived, at mainnet's settings: 3{,}935{,}930{,}637
      stroops; a keeper keeps the email contracts live for
      3{,}304{,}297{,}283 a year\claim{testnet-recovery-total}.
    \item Registry-root votes on a network: scheduled, not yet measured.
  \end{itemize}
\end{frame}

\begin{frame}{Evaluation: proving}
  \small
  \begin{tabular}{@{}lrr@{}}
    \toprule
    Prover (one desktop) & 18 threads & 1 thread \\
    \midrule
    Native \texttt{bb} & 0.94\,s & 4.64\,s \\
    bb.js in Node & 3.24\,s & 13.4\,s \\
    Chrome (wallet worker) & 1.44\,s & 8.37\,s \\
    Chrome, Performance-tier cores & 10.3\,s (4) & 32.5\,s \\
    Phone & \multicolumn{2}{c}{not measured} \\
    \bottomrule
  \end{tabular}
  \par\smallskip{\footnotesize\claims{desktop-proving-cost,browser-proving-cost,efficiency-core-proving}}
  \par\medskip
  \begin{itemize}\footnotesize
    \item macOS names the tiers Super (6) and Performance (12); load average
      about 6, a second run within 2\%.
    \item Chrome's renderer peaks at 1.4 to 2.5\,GiB; a phone's limits are
      untested. Budget of about two minutes: a target\claim{phone-proving-budget}.
  \end{itemize}
\end{frame}

\begin{frame}{Related work}
  \small
  \begin{itemize}
    \item zkEmail's recovery: email guardians from a hash of the address and
      an account code, a relayer that proves and stores the codes, and a
      registry with a delayed authorizer and per-user
      overrides~\cite{zkemail-recovery,zkemail-tx-builder,zkemail-user-overrideable-registry}.
    \item New here: the protocol's own statement, blinded; proofs in the
      browser; an immutable, attested registry with a salted, expiring pin;
      machine-checked composition; parsing checked on real mail.
    \item zkLogin and Aptos keyless make OIDC the key~\cite{zklogin,aptos-aip61}.
    \item Usable security: secret questions, trustees, recovery codes,
      login notifications and
      habituation~\cite{schechter-secret-questions,schechter-social-recovery,holtervennhoff-recovery-codes,markert-login-notifications,anderson-habituation}.
  \end{itemize}
\end{frame}

\begin{frame}{Limitations and open questions}
  \small
  \begin{itemize}
    \item Not built: the production enclave and its mail exchanger. All third
      parties are mocks. Not measured: phones, people. No audit; testnet only.
    \item The provider can link an email to its vote by timing; the sending
      service sees alert addresses.
    \item Open: pin semantics and lifetime, default actions, where the code
      is kept, registry parameters, the alert code's key, an on-chain
      enrollment proof, sender authentication per provider, who funds the
      keeper.
  \end{itemize}
\end{frame}

\begin{frame}{Summary}
  Email can be one recovery factor while the chain holds only a salted
  commitment: a browser-side DKIM proof over a blinded statement, and a
  per-account guardian inside Perch's delayed recovery. Soundness, replay and
  one-factor safety are proven over a model and privacy in an ideal model;
  the parsing and the model's agreement with the code are tested; every
  number on these slides carries its status.
\end{frame}

\begin{frame}[allowframebreaks]{References}
  \tiny
  \bibliographystyle{plain}
  \bibliography{references}
\end{frame}

\end{document}
